\subsection{ML in Nuclear Astrophysics}
% Source page: Deepfield.md Page 19; exactly one frame.
\setsource{}
\begin{frame}[c,label=md-19]

\sectiondivider{Reference}{Part II: ML in Nuclear Astrophysics}
\end{frame}

% Source page: Deepfield.md Page 34; exactly one frame.
\setsource{\href{https://doi.org/10.1103/PhysRevAccelBeams.25.044601}{Miskovich et al., PRAB 25, 044601 (2022)}}
\begin{frame}[c,label=md-34]{Exp Optimization: Online Bayesian Optimization}
\small
\begin{block}{Scientific question}
Can beam-based Bayesian optimization align SECAR and improve mass-plane
optics with fewer measurements than manual tuning?
\end{block}
\smallskip
\textbf{ML problem:} Sample-efficient optimization of measured beam metrics.
\vspace{0.3cm}

\slideheading{Mathematical Abstraction}
\medskip
\textbf{Alignment}
\[
\mathbf u_s^*=\arg\min_{\mathbf u_s}d_{\rm FP1}(\mathbf u_s)
\]
\smallskip
\textbf{Mass-plane optics}
\[
f_{\rm FP2}(\mathbf u_m)=
\left(\frac{x+w_x-x_0}{20}\right)^4+
\left(\frac{w_y}{70}\right)^4
\]
\slidenote{$\mathbf u_s,\mathbf u_m$: steerer and multipole settings;
$d_{\rm FP1}$: spot-center spread.\\[0.3em]
$x_0$: beam axis; $w_x,w_y$: spot widths.}
\end{frame}

% Source page: Deepfield.md Page 35; exactly one frame.
\setsource{\href{https://doi.org/10.1103/PhysRevAccelBeams.25.044601}{Miskovich et al., PRAB 25, 044601 (2022)}}
\begin{frame}[t,label=md-35]{Online Bayesian Optimization: Method and Results}
\small
\slideheading{Technical Stack}
\smallskip
\begin{itemize}
\setlength{\itemsep}{0.45em}
    \item \textbf{Diagnostics:} FP1 spot centers; FP2 position and widths.
    \item \textbf{Surrogate:} A Gaussian process, updated after each measurement.
    \item \textbf{Next setting:} Minimize $\mathrm{LCB}(\mathbf u)=\mu(\mathbf u)-\xi\sigma(\mathbf u)$.
    \item \textbf{Control:} GPy / GPyOpt + PyEpics; steerers, then multipoles.
\end{itemize}
\medskip
\slideheading{Results}
\smallskip
\begin{itemize}
\setlength{\itemsep}{0.4em}
    \item Angular deviation $<1$ mrad.
    \item At least \alert{$3\times$ faster} than manual tuning.
    \item \alert{32\%} smaller beam-spot width for $^{133}$Cs.
\end{itemize}
\end{frame}

% Source page: Deepfield.md Page 37; exactly one frame.
\setsource{\href{https://doi.org/10.1103/PhysRevResearch.7.013074}{Tsintari et al., PRR 7, 013074 (2025)}}
\begin{frame}[c,label=md-37]{Exp Optimization: SECAR $(p,n)$ Optics}
\small
\begin{block}{Scientific question}
Which SECAR settings separate nearly isobaric $(p,n)$ recoils from the beam
while retaining enough recoils for a cross-section measurement?
\end{block}
\smallskip
\textbf{ML problem:} Simulation-guided multi-objective optics optimization.
\vspace{0.3cm}

\slideheading{Mathematical Abstraction}
\medskip
\[
\mathbf u\in\mathbb R^{19}\xrightarrow{\text{COSY-INFINITY}}
(\Delta_{\rm FP2},\Delta_{\rm FP3},w_1,\ldots,w_K)
\]
\vspace{0.12cm}
\[
\begin{aligned}
\text{maximize}\quad & (\Delta_{\rm FP2},\Delta_{\rm FP3}) \\[0.8em]
\text{minimize}\quad & (w_1,\ldots,w_K)
\end{aligned}
\]
\slidenote[0.3cm]{%
$\Delta$: recoil--beam separation at slits;\quad
$w$: beam width along the separator.}
\end{frame}

% Source page: Deepfield.md Page 38; exactly one frame.
\setsource{\href{https://doi.org/10.1103/PhysRevResearch.7.013074}{Tsintari et al., PRR 7, 013074 (2025)}}
\begin{frame}[t,label=md-38]{SECAR $(p,n)$ Optics: Method and Results}
\small
\slideheading{Technical Stack}
\smallskip
\begin{itemize}
\setlength{\itemsep}{0.45em}
    \item \textbf{Search:} 19 field settings; recoil separation and beam widths.
    \item \textbf{Simulation:} COSY-INFINITY beam and recoil transport.
    \item \textbf{Optimization:} MOEA/D, Tchebycheff decomposition, differential evolution, and polynomial mutation.
    \item \textbf{Validation:} Beam tests of Pareto-selected optics; GP transmission interpolation.
\end{itemize}
\medskip
\slideheading{Results}
\smallskip

New ion optics enabled an astrophysical $(p,n)$ measurement:
\[
{}^{58}\mathrm{Fe}(p,n){}^{58}\mathrm{Co},\qquad
\boxed{\sigma=20.3\pm6.3\ \mathrm{mb}}.
\]
\end{frame}

